% ---------------------------------------------------------------------------
% Do the energies predict how far a sampled sequence's structure sits from the
% PTEN crystal? Subsampling of the equilibrated ladder, 2026-11-24.
%
% Every number here is measured and the file it comes from is named beside it.
% Sources: softmax/monte_carlo/energy_rmsd/{frames.tsv,scores.csv,
% summary_rmsd_xtal.txt,summary_rmsd_ref.txt,correlations_*.csv},
% softmax/monte_carlo/energy_rmsd_score.log, softmax/DEVLOG.txt 2026-11-24.
% Scripts: monte_carlo/energy_vs_rmsd.py, monte_carlo/plot_energy_rmsd.py.
% Companions: structure_comparison.tex (the 1D5R comparison itself),
% pten_k_sites_explained.md 6.6-6.6j (the two energies and why E is not usable
% in the loop), pten_metrics.tex (what a move costs).
%
% Build:  cd softmax/docs && pdflatex energy_rmsd.tex && pdflatex energy_rmsd.tex
% ---------------------------------------------------------------------------
\documentclass[11pt,a4paper]{article}

\usepackage[T1]{fontenc}
\usepackage[utf8]{inputenc}
\usepackage[margin=2.4cm]{geometry}
\usepackage{amsmath,amssymb}
\usepackage{booktabs}
\usepackage{graphicx}
\usepackage{array}
\newcolumntype{R}[1]{>{\raggedright\arraybackslash}p{#1}}
\usepackage[hidelinks]{hyperref}
\sloppy

\title{Do the energies know how far a sampled sequence is from the crystal?\\
\large Subsampling the equilibrated PTEN ladder, and what the answer rules out}
\author{Marco Sanchini}
\date{24 November 2026}

\begin{document}
\maketitle

\begin{abstract}
\noindent
Four thousand equilibrated sequences from the PTEN temperature ladder were decoded with
ESM3 and scored with every energy in the project, against the RMSD of each decode to the
1D5R crystal. The intended measurement did not survive contact with the data, and the
reason is the result. \textbf{ESM3 decodes every sampled sequence onto the PTEN fold.}
From 10 to 268 mutations out of 324 residues, the median CA RMSD to 1D5R moves from
2.01\,\AA{} to 2.30\,\AA{}, against 1.98\,\AA{} for the reference sequence itself; 95\% of
all four thousand frames sit below 2.42\,\AA. The target therefore has almost no range to
correlate with, and the decode's own schedule noise, 0.049\,\AA, is as large as the entire
interquartile spread of a cold chain. What correlation there is, is the mutation count:
pooled over the ladder every energy reaches $\rho\approx+0.63$ against the RMSD and the
Hamming distance alone reaches $+0.628$, and once that is partialled out $U_{am}$ retains
$+0.06$ and $U_{\mathrm{ce}}$ $+0.03$. The study is not wasted: against the RMSD to the
\emph{reference's own decode}, a coordinate with 4.6 times the spread and the one the
energy is actually defined against, the soft contact energy reaches $+0.95$ pooled and
$+0.60$ with the mutation count removed. A random-sequence control is the one measurement
needed before any of this can be written up as a property of the sampler rather than of
ESM3.
\end{abstract}

\section{What was done}

\subsection{The question, and the subsampling}

The sampler's energy $U_{am}$ is an attention-map distance computed from the reference
\emph{sequence}; no experimental structure enters it, and $U(\mathrm{ref})=0$ by
construction. A comparison against the crystal therefore tests \emph{interpretation} --
does a low energy mean the sequence folds like the real protein -- and not whether the
landscape is well posed. That is the question this study asks.

Four subsampling strategies were considered on 20 November and the dense one was chosen:
score a large, evenly spaced set of frames once on the GPU, then let the CPU pick any
subset from it, so no later selection costs GPU time. The frames are recovered without
re-running anything, because \texttt{data.dat} logs the site, the current residue, the
proposed residue and the accept flag of every move, so a chain's sequence at any move is
replayed from the reference on the CPU. Zero replay mismatches were found over all 4000
frames; the Hamming distance of every reconstructed sequence also matches the one the
sampler logged.

\begin{itemize}
\item \textbf{20 chains, 11 temperatures}: $T=1$, 1.5, 2, 2.5, 3, 4, 5, 6, 7, 8, 16. Two
chains each, except the hot-started $T=2.5$ and $T=7$ from urano, which ran one at a time.
\item \textbf{200 frames per chain, 4000 in all}, evenly spaced over the last 20\% of each
chain. The 0.8 burn-in is not the 0.5 default: at seven of the nine original temperatures
0.5 does not clear $t_{\mathrm{eq}}$.
\item \textbf{3976 of the 4000 sequences are distinct.} The 24 repeats are all at $T=1$,
where acceptance is 1.1\%; identical sequences are decoded once.
\end{itemize}

The two hot-started chains need their own treatment, since their $t_{\mathrm{eq}}$ is not
the ladder's relaxation time from the reference. Read off their own $U$ traces, $T=7$ is
flat from about 30\% of the chain onward and $T=2.5$ wanders by about $\pm80$ around a
spread of 30, which is what its autocorrelation time implies rather than a drift. The last
20\% is the same convention as the rest of the ladder, so it was kept.

\subsection{What was computed for each frame}

One ESM3 forward pass gives the attention map and the structure logits; one two-step
\texttt{generate()} plus the structure decoder gives coordinates. From those:

\begin{itemize}
\item $U_{am}$, taken from \texttt{data.dat} and recomputed from scratch for the first five
frames as a check (maximum difference: exactly 0).
\item $U_{\mathrm{ce}}^{\mathrm{ref}}$, the cross-entropy of the structure head against the
reference's own argmax structure tokens. This is the sampler's $U_{\mathrm{structure\_ce}}$,
the term that has had $\lambda=0$ for the whole project, computed in float32 to avoid the
bfloat16 quantisation trap of 18 November.
\item $U_{\mathrm{ce}}^{\mathrm{xtal}}$, the same cross-entropy against the VQ structure
tokens of 1D5R itself, over the 307 of 324 residues the crystal resolves.
\item $E^{\mathrm{ref}}$ and $E^{\mathrm{xtal}}$, the Zambon normalised contact-map energy
(CA, 8\,\AA, $|i-j|>2$) against the reference's decode at the same step count and against
the crystal's own contact map, with the soft-contact versions $E_s$ (sigmoid of width
1\,\AA{} in place of the hard cutoff) alongside at no extra cost.
\item The CA RMSD and the TM-score at the Kabsch superposition, to the crystal and to the
reference's decode; and the decode's mean pLDDT.
\item For 50 evenly spread frames, a second decode at 8 steps, for the noise floor.
\end{itemize}

\paragraph{Cost.} 15.4 minutes on astaroth, 0.231\,s per frame, peak 6276\,MB. The GPU was
never the constraint in this study; the statistics are.

\paragraph{The sanity checks all passed} (first lines of
\texttt{energy\_rmsd\_score.log}): $U_{\mathrm{ce}}^{\mathrm{ref}}(\mathrm{reference})
= 2.397402$, matching the value measured independently on 18 November; the same sequence
decoded twice differs in 0 contacts, so the decode is deterministic at fixed step count;
re-encoding the reference's own decode recovers 37.3\% of the tokens it was decoded from,
far above the $1/4096 = 0.02\%$ a broken encoder call would give.

\subsection{How many independent frames there really are}

Each chain's integrated autocorrelation time of $U$ was measured inside its own window.
It runs from 1264 moves ($T=6$) to 11005 ($T=2$), which leaves between 10 and 58
independent frames per chain and about 404 over the whole ladder. The 4000 frames are
therefore oversampled on purpose: the extra frames cost nothing and make the CPU-side
subsets possible, but every interval below is a moving-block bootstrap with blocks one
autocorrelation time long, never a naive $n$.

\section{Results}

\subsection{The headline: the decode never leaves the fold}

\begin{table}[h]
\centering
\begin{tabular}{lrrrrrrr}
\toprule
$T$ & frames & mutations & \multicolumn{3}{c}{RMSD to 1D5R [\AA]} & TM & pLDDT \\
\cmidrule(lr){4-6}
 & & & min & median & max & median & median \\
\midrule
reference & 1 & 0 & & 1.98 & & 0.942 & 88.8 \\
\midrule
1    & 400 & 10--28   & 1.93 & 2.01 & 2.11 & 0.940 & 88.2 \\
1.5  & 400 & 35--59   & 1.88 & 2.00 & 2.14 & 0.941 & 87.9 \\
2    & 400 & 54--92   & 1.88 & 2.00 & 2.25 & 0.941 & 88.0 \\
2.5  & 200 & 87--106  & 1.90 & 2.02 & 2.16 & 0.939 & 87.8 \\
3    & 400 & 115--150 & 1.84 & 2.02 & 2.22 & 0.939 & 87.3 \\
4    & 400 & 144--191 & 1.88 & 2.05 & 2.60 & 0.937 & 86.6 \\
5    & 400 & 170--212 & 1.89 & 2.08 & 3.03 & 0.935 & 86.1 \\
6    & 400 & 193--222 & 1.87 & 2.08 & 2.49 & 0.935 & 85.9 \\
7    & 200 & 201--230 & 1.93 & 2.17 & 3.60 & 0.927 & 85.4 \\
8    & 400 & 207--239 & 1.90 & 2.16 & 3.42 & 0.928 & 85.2 \\
16   & 400 & 232--268 & 1.98 & 2.30 & 4.00 & 0.918 & 83.2 \\
\bottomrule
\end{tabular}
\caption{Every temperature of the ladder, 4000 frames. At $T=16$ the sequences differ from
the reference at 232 to 268 of 324 positions -- up to 83\% of the protein -- and ESM3 still
places them 2.30\,\AA{} from the crystal, with a TM-score of 0.92. Only 3.0\% of all frames
exceed 2.5\,\AA{} and 14 of 4000 exceed 3\,\AA.}
\end{table}

This is the result of the study, and it was not the one being looked for. The target that
the whole design was built around has an interquartile spread of 0.148\,\AA{} over the
entire ladder, and between 0.046\,\AA{} ($T=1$) and 0.249\,\AA{} ($T=16$) within a chain.
Panel~A of Figure~\ref{fig:main} is the picture: a flat band just above the reference's own
1.98\,\AA, with a thin tail appearing only past 200 mutations.

Two readings are possible and they are not yet separated. Either the sampler genuinely
preserves the fold at every temperature -- which is what one would like to conclude -- or
ESM3's structure head returns something close to the PTEN fold for a wide class of
sequences, in which case the flat band says nothing about the sampler at all. Section~3.1
says how to tell them apart, and until that is done the honest statement is the second one.

\begin{figure}[t]
\centering
\includegraphics[width=\textwidth]{../monte_carlo/energy_rmsd/energy_rmsd_fig.png}
\caption{\textbf{A}: the RMSD to 1D5R against the number of mutations, all 4000 frames,
coloured by sampling temperature. \textbf{B}: the mean pLDDT of the same decodes, which
does fall, smoothly and with real range. \textbf{C}: the soft contact energy against the
RMSD to the \emph{reference's} decode, the coordinate that has range. \textbf{D}: 50 frames
decoded at 2 and at 8 steps -- the scatter about the diagonal is the measurement noise of
panel~A's vertical axis.}
\label{fig:main}
\end{figure}

\subsection{Pooled over the ladder, the energies know the mutation count and little else}

\begin{table}[h]
\centering
\begin{tabular}{lcccc}
\toprule
 & \multicolumn{2}{c}{vs.\ RMSD to 1D5R} & \multicolumn{2}{c}{vs.\ RMSD to the reference's decode} \\
\cmidrule(lr){2-3}\cmidrule(lr){4-5}
energy & $\rho$ & $\rho$ given $H_d$ & $\rho$ & $\rho$ given $H_d$ \\
\midrule
$U_{am}$                      & $+0.629$ & $+0.061$ & $+0.927$ & $+0.189$ \\
$U_{\mathrm{ce}}^{\mathrm{ref}}$ & $+0.625$ & $+0.030$ & $+0.923$ & $+0.139$ \\
$E^{\mathrm{ref}}$            & $+0.641$ & $+0.176$ & $+0.939$ & $+0.481$ \\
$E_s^{\mathrm{ref}}$          & $+0.652$ & $+0.229$ & $\mathbf{+0.952}$ & $\mathbf{+0.597}$ \\
\midrule
$U_{\mathrm{ce}}^{\mathrm{xtal}}$\,* & $+0.622$ & $+0.114$ & $+0.883$ & $+0.092$ \\
$E^{\mathrm{xtal}}$\,*        & $+0.655$ & $+0.304$ & $+0.804$ & $+0.148$ \\
$E_s^{\mathrm{xtal}}$\,*      & $+0.692$ & $+0.377$ & $+0.889$ & $+0.214$ \\
\midrule
$H_d$ itself                  & $+0.628$ & --- & $+0.925$ & --- \\
\bottomrule
\end{tabular}
\caption{Spearman correlations over all 4000 frames; every interval (in
\texttt{summary\_rmsd\_xtal.txt}) is a block bootstrap with blocks one autocorrelation time
long and is within $\pm0.03$ of the estimate. Rows marked * are energies that look at the
crystal themselves, so part of their correlation with a crystal-based target is built in by
construction. The 404-frame decorrelated subset and the 258-frame energy-stratified subset
give the same picture to within 0.07.}
\end{table}

Read the first column against the last row. Every energy lands within 0.03 of the Hamming
distance's own $+0.628$, and the partial correlations say why: with the mutation count
removed, $U_{am}$ keeps $+0.06$ and the structure cross-entropy $+0.03$. Neither of the two
energies the sampler can actually use knows anything about the crystal beyond how many
mutations have accumulated.

\paragraph{Three things do survive the partialling, and two of them are circular.}
$E^{\mathrm{xtal}}$ and $E_s^{\mathrm{xtal}}$ reach $+0.30$ and $+0.38$, but they are read
off the same decode as the RMSD and against the same crystal, so this is partly an identity
rather than a finding -- it was flagged as such before the run. The honest survivor is
$E_s^{\mathrm{ref}}$ at $+0.23$, the soft contact energy against the reference's prediction,
which never sees the crystal.

\paragraph{The two in-loop energies have collapsed onto each other.} On these frames
$U_{am}$ and $U_{\mathrm{ce}}^{\mathrm{ref}}$ agree at $\rho=+0.992$. That is a change of
regime worth recording: on the 324$\times$19 single-mutant scan of 18 November they were
complementary, the cross-entropy winning per site and $U_{am}$ per substitution. At 50 to
270 mutations both have become monotone functions of distance from the reference, and
$U_{\mathrm{ce}}$ adds nothing to $U_{am}$ as a rescoring term at this scale. This also
disposes, for the moment, of the question of what $\lambda_{\mathrm{structure\_ce}}$ should
be: a term that correlates at $+0.99$ with the one already in the energy will not change
the sampling, whatever its scale.

\subsection{A coordinate that does have range}

The right-hand half of the table is the same analysis against the RMSD to the reference's
\emph{own} decode. That coordinate spans 0.29 to 3.47\,\AA{} with an interquartile range of
0.68\,\AA, 4.6 times the crystal RMSD's, and it is the structurally natural question to ask
of this energy, which is defined against the reference and not against any experiment.
Against it the soft contact energy reaches $+0.95$ pooled and $+0.60$ with the mutation
count removed (Figure~\ref{fig:main}C), and the hard $E^{\mathrm{ref}}$ $+0.48$. Matched
across temperatures, $E_s^{\mathrm{ref}}$ holds between $+0.44$ and $+0.77$ with no trend
in $T$, so this is a property of the energy and not of one corner of the ladder.

The mean pLDDT behaves similarly: it falls smoothly from 88.2 to 83.2 along the ladder
(Figure~\ref{fig:main}B), and $E_s^{\mathrm{ref}}$ predicts it at $-0.34$ with the mutation
count removed, against $-0.11$ for $U_{am}$ (negative because a high pLDDT is a good
structure, so the informative sign is reversed for this target).

\subsection{The measurement noise of the target}

The 50 frames decoded a second time at 8 steps instead of 2 shift by a median of
0.049\,\AA{} in their RMSD to the crystal, and the two step counts agree only at
$\rho=+0.65$. Compare that with the within-chain spread the measurement has to resolve:
0.046\,\AA{} at $T=1$, 0.249\,\AA{} at $T=16$. At the cold end the schedule noise is larger
than the whole signal; at the hot end it is a fifth of it. The contact-map energies are far
better behaved under the same test ($\rho=+0.96$ for $E^{\mathrm{ref}}$, a median shift of
0.11 of its interquartile range), which is consistent with the September--November finding
that $E$ is usable per site but not per mutant.

This is a second, independent reason not to build anything on the crystal RMSD as measured
here, and it is a property of the decoder's schedule, not of the sampler.

\subsection{Per temperature}

Within a single temperature the crystal RMSD is essentially unpredictable at the cold end
and mildly predictable at the hot end: for $E_s^{\mathrm{ref}}$, $+0.01$ at $T=1$ rising to
$+0.60$ at $T=7$ and $+0.25$ at $T=16$, with $U_{am}$ never leaving the band $-0.12$ to
$+0.30$. These numbers are taken from the matched table, in which every chain is cut to the
same number of its own autocorrelation times, because the raw per-temperature row is
biased by how much of its equilibrium spread each window happens to cover: a chain that
wanders less gives a smaller $|\rho|$ for that reason alone. This is range restriction, and
it is a different mechanism from the window attenuation that produced the spurious $T=6$
entropy peak, but it has the same consequence -- a raw temperature trend cannot be trusted.
Both tables, and the per-temperature spreads that justify the correction, are printed by
\texttt{energy\_vs\_rmsd.py analyze}.

\section{What needs to be done}

\subsection{The one measurement that decides what Section 2.1 means}

\textbf{Decode random sequences of the same length.} If a random 324-residue sequence also
lands 2.3\,\AA{} from 1D5R with pLDDT in the eighties, then the flat band is a property of
ESM3's structure head under a sequence-only prompt and says nothing whatever about the
sampler, and every structural statement in this document has to be rewritten as a statement
about the model. If random sequences instead land far away, then the sampler really does
keep the fold at 83\% sequence divergence, which is a strong and publishable claim about
the energy. Nothing else in this programme should be built until this is known.
\texttt{bench\_decode.py} already has a \texttt{--random} flag and the machinery; ten random
sequences and ten shuffles of the reference (same composition, no fold) cost under a minute
on astaroth.

\subsection{Then, in order}

\begin{enumerate}
\item \textbf{Move the target to the reference's decode.} If the control comes back
favourable, the RMSD to the reference's own prediction, and the mean pLDDT, are the
coordinates with range and with a clean interpretation for an energy defined against the
reference. The analysis already runs against both, \texttt{--target rmsd\_ref} and a
pLDDT target being one flag apart.

\item \textbf{Settle whether the soft contact energy is adopted.} It is the only
non-circular energy with signal beyond the mutation count here ($+0.60$ against the
reference RMSD), and it costs nothing extra once a decode exists. Against that, it is a
departure from the definition in Zambon et al, so it belongs in the thesis as a diagnostic
unless there is a reason to redefine $E$. This is a decision for the supervisors, not a
measurement.

\item \textbf{Re-run when the chains grow.} The sampler is being extended. Regenerating
\texttt{frames.tsv} and rescoring costs about a quarter of a second per frame, so a full
redo is 15--20 minutes; \texttt{analyze} drops any scored row that is not in the current
frame list, so a resumed \texttt{scores.csv} cannot silently mix two windows. The cheapest
real improvement is more chains rather than longer ones: 404 independent frames across 11
temperatures is what limits every per-temperature number above.

\item \textbf{If the crystal RMSD is wanted as a serious observable,} it needs a
deterministic single-pass structure predictor rather than a two-step \texttt{generate()}.
ESMFold is the obvious candidate and is what Zambon et al use. The alternative, averaging
several schedules, was measured in November not to work: the mean rank of five step counts
is no better than the best single one.

\item \textbf{Leave $\lambda_{\mathrm{structure\_ce}}$ at zero for now.} The scale
calculation of 18 November suggested $\lambda\sim10^3$ would be needed for the term to
matter at all, and Section~2.2 shows that at the distances these chains explore the term is
$+0.99$ correlated with $U_{am}$ anyway. There is no case for switching it on until it has
been shown to disagree with $U_{am}$ somewhere that matters.
\end{enumerate}

\subsection{Open questions worth one line each}

\begin{itemize}
\item The crystal's VQ structure tokens agree with the reference's argmax tokens at only
12.1\% of resolved residues, yet the two structures are 1.98\,\AA{} apart. The structure
token alphabet is evidently far finer than the geometry it encodes; this should be checked
before $U_{\mathrm{ce}}^{\mathrm{xtal}}$ is quoted anywhere as a structural distance.
\item 17 of 324 residues are unresolved in 1D5R and are excluded throughout. They are not
randomly placed, and whether any of them carry the sampler's interesting sites has not been
checked.
\item The single worst frame, \texttt{T16.c1@176482} at 4.00\,\AA{} and 256 mutations, still
has pLDDT 81. Whatever the decoder is doing with these sequences, it is not reporting
uncertainty about it.
\end{itemize}

\end{document}
